\documentclass[11pt]{article}

\usepackage{hyperref}
\usepackage[letterpaper, top=1in, bottom=1in, left=1.5in, right=1in]{geometry}
\usepackage{graphicx}
\usepackage{wrapfig}
\usepackage{epstopdf}
\usepackage{pdfpages}
\usepackage[titletoc,title]{appendix}
\usepackage[nottoc,notbib]{tocbibind}
\usepackage{lipsum}
\usepackage[toc]{glossaries}
\usepackage{verbatim}
\usepackage{subcaption}
\usepackage{bookmark}
\usepackage{setspace}
\usepackage{titlesec}
\usepackage{setspace}

% Font
% Title and Heading Styling
%%%%%%%%%%%%%%%%%%%%%%%%%%%%%%%%%%%%%%%%%%%%%%%%%%%%%%%%%%%%%%%%%%%%%%%%%%%%%%

\definecolor{heading-color}{HTML}{1A8EBD}
\definecolor{subheading-color}{HTML}{333333}
\definecolor{subsubheading-color}{HTML}{666666}

\titleformat{\section}
  {\normalfont\sffamily\fontsize{16pt}{16}\bfseries\color{heading-color}}
  {\thesection}{1em}{}

\titleformat{\subsection}
  {\normalfont\sffamily\fontsize{14pt}{14}\bfseries\color{subheading-color}}
  {\thesubsection}{1em}{}
\titleformat{\subsubsection}
  {\normalfont\sffamily\fontsize{14pt}{14}\bfseries\color{subsubheading-color}}
  {\thesubsubsection}{1em}{}

% Section Tweaks
%%%%%%%%%%%%%%%%%%%%%%%%%%%%%%%%%%%%%%%%%%%%%%%%%%%%%%%%%%%%%%%%%%%%%%%%%%%%%%

\renewcommand{\abstractname}{Executive Summary}

\usepackage{etoolbox}
\patchcmd{\thebibliography}{\section*{\refname}}{}{}{}

\bookmarksetup{open,numbered}


%   Glossary
%%%%%%%%%%%%%%%%%%%%%%%%%%%%%%%%%%%%%%%%%%%%%%%%%%%%%%%%%%%%%%%%%%%%%%%%%%%%%%


\newglossaryentry{philipsHue}
{
  name=Philips Hue Light Bulb,
  description={A ``smart'' light bulb that connects
  to home Wi-Fi networks, \url{http://meethue.com/}}
}

\newacronym{api}{API}{Application Programming Interface}

\makeglossaries


%   _____                _     __  __       _   _            
%  |  ___| __ ___  _ __ | |_  |  \/  | __ _| |_| |_ ___ _ __ 
%  | |_ | '__/ _ \| '_ \| __| | |\/| |/ _` | __| __/ _ \ '__|
%  |  _|| | | (_) | | | | |_  | |  | | (_| | |_| ||  __/ |   
%  |_|  |_|  \___/|_| |_|\__| |_|  |_|\__,_|\__|\__\___|_|   
%



\begin{document}

\raggedright

\pagenumbering{gobble}

%   Letter of Transmittal
%%%%%%%%%%%%%%%%%%%%%%%%%%%%%%%%%%%%%%%%%%%%%%%%%%%%%%%%%%%%%%%%%%%%%%%%%%%%%%

    %\pdfbookmark[section]{Letter of Transmittal}{lot}
    %\includepdf{../"Letter of Transmittal"/"Letter of Transmittal".pdf}
   

%   Cover Page
%%%%%%%%%%%%%%%%%%%%%%%%%%%%%%%%%%%%%%%%%%%%%%%%%%%%%%%%%%%%%%%%%%%%%%%%%%%%%%
    
    \pdfbookmark[section]{Cover Page}{cover}
    \includepdf{Pages/CoverPage.pdf}


%   Table of Contents and Figures 
%%%%%%%%%%%%%%%%%%%%%%%%%%%%%%%%%%%%%%%%%%%%%%%%%%%%%%%%%%%%%%%%%%%%%%%%%%%%%%

    \setcounter{page}{1}
    \pagenumbering{roman}
    

\setlength{\parskip}{1mm}   
    
    \pdfbookmark[section]{\contentsname}{toc}
    \tableofcontents
    \newpage
    

\setlength{\parskip}{2mm}   
    
    
    \phantomsection
    \listoffigures
    \phantomsection
    \listoftables
%    \newpage
    
    \printglossaries
    \newpage
        
%   Summary 
%%%%%%%%%%%%%%%%%%%%%%%%%%%%%%%%%%%%%%%%%%%%%%%%%%%%%%%%%%%%%%%%%%%%%%%%%%%%%%%
    
\begin{abstract}
\phantomsection
\addcontentsline{toc}{section}{Executive Summary}
\noindent  

Good Morning is an app backed by \textit{SoftStart}. The app will help users wake up in the morning using a personal assistant, machine learning and smart integrations. Despite a number of other apps in this domain we believe that \textbf{Good Morning} will provide users with a wake up experience currently not available in other applications. The objectives for \textbf{Good Morning} are,

\begin{description}
  \item[An Alarm] Core alarm functionality or similar to wake the user up gracefully.
  \item[Smart Assistant] A personal assistant to help take you from waking up, through your morning routine and out the door.
  \item[Machine Learning] Support for learning the users morning routine and assisting them in adjusting their schedule and wake up time accordingly.
  \item[Smart Integration] Connection to smart devices around the house (e.g. a \gls{philipsHue}) to help the user throughout their morning routine.
\end{description}

\textit{SoftStart} is excited to start developing the requirements for \textbf{Good Morning} with your organization. We look forward to hearing from you!

\end{abstract}
\newpage

%   ____            _       
%  | __ )  ___   __| |_   _ 
%  |  _ \ / _ \ / _` | | | |
%  | |_) | (_) | (_| | |_| |
%  |____/ \___/ \__,_|\__, |
%                     |___/ 
    
    
\setcounter{page}{1}
\pagenumbering{arabic}
\glsresetall
\renewcommand{\arraystretch}{1.5} % Fix the table sizes
    
%%%%%%%%%%%%%%%%%%%%%%%%%%%%%%%%%%%%%%%%%%%%%%%%%%%%%%%%%%%%%%%%%%%%%%%%%%%%%%%%
    
%
% INTRODUCTION
% ---------------------------------------
\section{Introduction}\label{introduction}

Here at \textit{SoftStart} software incubator we make great ideas happen.
\textbf{Good Morning} is an app that will help users wake up in the
morning. It will help them get the most out of their morning routine with a personal assistant. In addition it will use machine learning and smart integrations to make their morning the best it can be. More information about the problem and objective can be found in sections \ref{problem-description} and \ref{project-objectives} respectively.

% TODO: BEN - Current Systems intro

\textbf{Good Morning} will focus on the young adults and young professionals. Users who are interested in increasing their productivity and organization and are familiar with technology for day-to-day use (see section \ref{target-users}). Many of these users may already have smart, internet of things, devices in the household. \textbf{Good Morning} will integrate with these devices to provide further support during the morning routine. Possible integrations are discussed further in section \ref{known-interactions}.

A significant constraint for \textbf{Good Morning} will be internationalization. Because of the integrated nature of \textbf{Good Morning} there will be issues around cultural norms as well as the personal assistant. See section \ref{known-constraints} for more information on constraints.

%
% PROBLEM DEFINITION
% ---------------------------------------
\section{Problem Description}\label{problem-description}



% TODO: ANDREI - Problem Definition

%
% PROJECT OBJECTIVES
% ---------------------------------------
\section{Project Objectives}\label{project-objectives}

There are several key objectives that make up \textbf{Good Morning},

\begin{itemize}
\item
  Deliver a system that reliably wakes the user up.
\item
  Deliver additional information to improve the process of waking up, such as Weather, Missed Messages etc.
\item
  A flexible system that can adapt to changes in your schedule.
\item
  Integration with additional smart technologies, e.g. a \glslink{philipsHue}{Philips Hue wakeup light}
\end{itemize}

% TODO: JAKE - Project Objectives

%
% CURRENT SYSTEMS
% ---------------------------------------
\section{Current Systems}\label{current-systems}

Softstart believes that \textbf{Good Morning} has the potential to fill
a sizeable market niche that is currently vacant. There are currently no other all-in-one systems
for boosting morning productivity. There are however some other applications
on the market that provide similar tools and functions.

Wake Up Applications: There are many different wake up applications on the market, most
of which are simple alarm clocks, the following are a few of the more popular and unique ones in use.
\begin{itemize}
\item
    \textbf{Sleep Cycle}: Sleep cycle is a simple alarm clock that monitors your sleep
    cycle, and wakes you up during the lightest part of your sleep.
\item
    \textbf{iCuckoo}: This alarm application makes a donation to a charity of your choice every time you hit the
    the snooze button.
\item
    \textbf{Wakie}: Wakie is a different but unique approach to waking up. It
    involves two anonymous users. Users are either wakies or sleepyheads. 
    The sleepyheads determine what time they would like to wake up and how they would
    like to be woken up. Then the wakies will contact them at the specified time and 
    discuss whatever topic the sleepyhead specified. The wakies only have 1 minute to
\item
	\textbf{Google Now}: Google Now provides useful information based on your current location, frequently
	visited locations, and search history. It provides a lot of useful and relevant information to users
	depending on the time of day.
\end{itemize}

\textbf{Good Morning} will provide users with many of the tools that the above applications
provide, but package them all together with an easy to use interface. On top of that, it will also
interact with other systems through the internet of things, allowing users to have even more flexibility
and utility. No other applications have achieved this yet

% TODO: JONAH - Current Systems

%
% TARGET USERS
% ---------------------------------------
\section{Target Users}\label{target-users}

Good Morning targets the age group of 18 to 32 because this age group is statistically the most likely to suffer from sleep issues \cite{sleepdisorders}. Furthermore, this age group is extremely busy and high stress, and Good Morning can help by cognitively offloading all of the stress of morning planning. This group of users can benefit greatly from Good Morning technology to improve their morning routine, which affects the rest of their day.

By utilizing the Good Morning application, this demographic of users will be able to improve all aspects of their life by waking up pleasantly and comfortably.
% TODO: JAKE - Target Users

%
% KNOWN INTERACTIONS
% ---------------------------------------
\section{Known Interactions}\label{known-interactions}

\begin{itemize}
\item
  Smart devices, e.g. \gls{philipsHue}
\item
  Mobile Operating System and \glspl{api}, i.e., Android or iOS
\end{itemize}

% TODO: ANDREI - Known Interactions

%
% Known Contraints
% ---------------------------------------
\section{Known Constraints}\label{known-constraints}

\begin{itemize}
 \item The app must be available on various mobile platforms
 \item Project should be completed by the end of April
 \item Must be available and useable internationally and support multilanguage functionality
 
\end{itemize}

% TODO: JONAH - Constaints

%
% Project Schedule
% ---------------------------------------
\section{Project Schedule}\label{project-schedule}

The requirements for \textbf{Good Morning} will be developed according to a rapid iteration schedule with the final prototype being completed on March 31st, 2016. The full project schedule is show in table \ref{tab:schedule}.

\begin{table}[htb]
    \caption{\textbf{Good Morning} Project Schedule}\label{tab:schedule}
    \centering
    \begin{tabular}{|lr|}
        \hline
        \textbf{Project Milestone} & \textbf{Date} \\
        \hline
        Draft Requirements Definition & 21 Jan \\
        Promotional Website Live & 21 Jan \\
        Formal Requirements Definition (Revision) & 16 Feb \\
        Detailed Requirements Definition (Revision) & 1 Mar \\
        Preliminary Prototype Demo & 3 Mar \\
        Final Requirements Definition (Revision) & 15 Mar \\
        User Manual & 22 Mar \\
        Final Prototype Demo & 31 Mar \\
        \hline
    \end{tabular}
\end{table}



%
% PROJECT TEAM
% ---------------------------------------
\section{Project team}\label{project-team}

The development of the requirements for \textbf{Good Morning} will require a variety of skill sets across several disciplines. The skill sets required will be

\begin{description}
  \item[Product Manager] to plan and organize the development of prototypes and specifications.
  \item[Writer] to write and maintain project documentation and copy.
  \item[Web Developer] to plan and maintain the promotional website.
  \item[Mobile Developer(s)] to plan and develop requirements and prototypes.
  \item[Designer] to plan and develop user driven requirement and prototype user-experience options.
  \item[Market Analyst] to develop market based requirements.
\end{description}



    
%%%%%%%%%%%%%%%%%%%%%%%%%%%%%%%%%%%%%%%%%%%%%%%%%%%%%%%%%%%%%%%%%%%%%%%%%%%%%%%
 
    \pagebreak
    \section{References} 
    
    \bibliographystyle{IEEEtran}

    \bibliography{references}
    
    \pagebreak
%   ____             _      __  __       _   _            
%  | __ )  __ _  ___| | __ |  \/  | __ _| |_| |_ ___ _ __ 
%  |  _ \ / _` |/ __| |/ / | |\/| |/ _` | __| __/ _ \ '__|
%  | |_) | (_| | (__|   <  | |  | | (_| | |_| ||  __/ |   
%  |____/ \__,_|\___|_|\_\ |_|  |_|\__,_|\__|\__\___|_|   
%  

% Appendices

%\pdfbookmark[part]{Appendices}{appx}
%\begin{appendices}
%
%\section{Tool Screen Captures}\label{appx:screens}
%
%
%\end{appendices}

\end{document}
